
% ==========================================================

\ifnum\llncs=0\ifnum\fullversion=1\ifnum\anonymous=0
\renewcommand*{\backref}[1]{(Cited on page~#1.)}
\fi\fi\fi

\DeclareMathAlphabet{\mathsl}{OT1}{cmr}{m}{sl}
\DeclareMathAlphabet{\mathsc}{OT1}{cmr}{m}{sc}
\DeclareMathAlphabet{\mathslbf}{OT1}{cmr}{bx}{sl}
% math script font; extra commands to make slightly larger
\DeclareFontFamily{OT1}{pzc}{}
\DeclareFontShape{OT1}{pzc}{m}{it}%
             {<-> s * [1.150] pzcmi7t}{}
\DeclareMathAlphabet{\mathscript}{OT1}{pzc}{m}{it}
\setlength{\fboxsep}{1pt}

% ====================================================================

% Theorem environments

\ifnum\llncs=0

\newtheorem{thm}{Theorem}[section]
\newtheorem{lem}[thm]{Lemma}
\newtheorem{cor}[thm]{Corollary}
\newtheorem{propo}[thm]{Proposition}
\newtheorem{clm}[thm]{Claim}
\newtheorem{defn}[thm]{Definition}
\newtheorem{assm}[thm]{Assumption}
\newtheorem{rem}[thm]{Remark}
\newtheorem{obs}[thm]{Observation}
\newtheorem{egs}[thm]{Example}
\newtheorem{fct}[thm]{Fact}
\newtheorem{expr}{Experiment}
\newtheorem{cons}[thm]{Construction}
\newtheorem{nte}[thm]{Note}
\newtheorem{conj}[thm]{Conjecture}
\newtheorem{task}[thm]{Task}
\newtheorem{ques}[thm]{Question}
\newtheorem{idea}[thm]{New Idea}
%\newtheorem{task}[thm]{Task}

\newenvironment{theorem}{\begin{thm}}{\end{thm}}
\newenvironment{lemma}{\begin{lem}}{\end{lem}}
\newenvironment{corollary}{\begin{cor}}{\end{cor}}
\newenvironment{proposition}{\begin{propo}}{\end{propo}}
\newenvironment{definition}{\begin{defn}}%
{\end{defn}}
%\newenvironment{definition}{\begin{defn}\begin{rm}}%
%{\end{rm}\end{defn}}
\newenvironment{assumption}{\begin{assm}\begin{em}}%
{\end{em}\end{assm}}
\newenvironment{claim}{\begin{clm}\begin{rm}}%
{\end{rm}\end{clm}}
\newenvironment{remark}{\begin{rem}\begin{em}}%
{\end{em}\end{rem}}
\newenvironment{observation}{\begin{obs}\begin{em}}%
{\end{em}\end{obs}}
\newenvironment{example}{\begin{egs}\begin{em}}%
{\end{em}\end{egs}}
\newenvironment{fact}{\begin{fct}\begin{em}}%
{\end{em}\end{fct}}
\newenvironment{construction}{\begin{cons}\begin{rm}}%
{\end{rm}\end{cons}}
\newenvironment{note}{\begin{nte}\begin{rm}}%
{\end{rm}\end{nte}}

\else

\ifnum\fullversion=1
\newtheorem{thm}{Theorem}
\newtheorem{lem}[thm]{Lemma}
\newtheorem{cor}[thm]{Corollary}
\newtheorem{propo}[thm]{Proposition}
\newtheorem{clm}[thm]{Claim}
\newtheorem{defn}[thm]{Definition}
\newtheorem{assm}[thm]{Assumption}
\newtheorem{rem}[thm]{Remark}
\newtheorem{fct}[thm]{Fact}
\newtheorem{expr}{Experiment}
\newtheorem{cons}[thm]{Construction}

\renewenvironment{theorem}{\begin{thm}}{\end{thm}}
\renewenvironment{lemma}{\begin{lem}}{\end{lem}}
\renewenvironment{corollary}{\begin{cor}\begin{rm}}%
{\end{rm}\end{cor}}
\renewenvironment{proposition}{\begin{propo}\begin{rm}}%
{\end{rm}\end{propo}}
\renewenvironment{definition}{\begin{defn}\begin{rm}}%
{\end{rm}\end{defn}}
\newenvironment{assumption}{\begin{assm}\begin{em}}%
{\end{em}\end{assm}}
\renewenvironment{claim}{\begin{clm}\begin{rm}}%
{\end{rm}\end{clm}}
\renewenvironment{remark}{\begin{rem}\begin{em}}%
{\end{em}\end{rem}}
\newenvironment{fact}{\begin{fct}\begin{em}}%
{\end{em}\end{fct}}
\newenvironment{construction}{\begin{cons}\begin{rm}}%
{\end{rm}\end{cons}}
\fi

\fi


\newenvironment{sfigure}{\begin{figure}[t]\begin{small}}{\end{small}\end{figure}}

\newcommand{\secref}[1]{Section~\ref{#1}}
\newcommand{\apref}[1]{Appendix~\ref{#1}}
\newcommand{\thref}[1]{Theorem~\ref{#1}}
\newcommand{\defref}[1]{Definition~\ref{#1}}
\newcommand{\corref}[1]{Corollary~\ref{#1}}
\newcommand{\lemref}[1]{Lemma~\ref{#1}}
\newcommand{\clref}[1]{Claim~\ref{#1}}
\newcommand{\propref}[1]{Proposition~\ref{#1}}
\newcommand{\factref}[1]{Fact~\ref{#1}}
\newcommand{\remref}[1]{Remark~\ref{#1}}
\newcommand{\obsref}[1]{Observation~\ref{#1}}
\ifnum\llncs=1
  \newcommand{\figref}[1]{Fig.~\ref{#1}}
\else
  \newcommand{\figref}[1]{Figure~\ref{#1}}
\fi
\newcommand{\tabref}[1]{Table~\ref{#1}}
\newcommand{\consfigref}[1]{Construction~\ref{#1}}
\renewcommand{\eqref}[1]{\mbox{Equation~(\ref{#1})}}
\newcommand{\Eqref}[1]{\mbox{Eq.~(\ref{#1})}}
\newcommand{\Eqsref}[2]{Eqs.~(\ref{#1}) and~(\ref{#2})}
\newcommand{\consref}[1]{Construction~\ref{#1}}


\def\sqed{{\hspace{1pt}\rule[-1pt]{3pt}{9pt}}}
\def\qedsym{\hspace{1pt}\rule[-1pt]{3pt}{9pt}}
\newcommand{\mybox}{\raisebox{-1.5pt}{{\Large $\Box$}}}
\def\enddef{{\sqed}}

\newlength{\saveparindent}
\setlength{\saveparindent}{\parindent}
\newlength{\saveparskip}
\setlength{\saveparskip}{\parskip}

\ifnum\llncs=0
% deluxe proof environment
\def\qsym{\vrule width0.6ex height1em depth0ex}
\newcount\proofqeded
\newcount\proofended
\def\qed{{\hspace{1pt}\rule[-1pt]{3pt}{9pt}}
\end{rm}\addtolength{\parskip}{-0pt}
\setlength{\parindent}{\saveparindent}
\global\advance\proofqeded by 1 }
\def\qedenv{
\end{rm}\addtolength{\parskip}{-0pt}
\setlength{\parindent}{\saveparindent}
\global\advance\proofqeded by 1 }
\newenvironment{proof}%
 {\proofstart}%
 {\ifnum\proofqeded=\proofended~\qed\fi \global\advance\proofended by 1
  \medskip}
\newenvironment{proofenv}%
 {\proofenvstart}%
 {\ifnum\proofqeded=\proofended\qedenv\fi \global\advance\proofended by 1
  \medskip}
\makeatletter
\def\proofstart{\@ifnextchar[{\@oprf}{\@nprf}}
\def\proofenvstart{\@ifnextchar[{\@osprf}{\@nsprf}}
\def\@oprf[#1]{\begin{rm}\protect\vspace{6pt}\noindent{\bf Proof of #1:\ }%
\addtolength{\parskip}{5pt}\setlength{\parindent}{0pt}}
\def\@osprf[#1]{\begin{rm}\protect\vspace{6pt}\noindent
\addtolength{\parskip}{5pt}\setlength{\parindent}{0pt}}
\def\@nprf{\begin{rm}\protect\vspace{6pt}\noindent{\bf Proof:\ }%
\addtolength{\parskip}{5pt}\setlength{\parindent}{0pt}}
\def\@nsprf{\begin{rm}\protect\vspace{6pt}\noindent%
\addtolength{\parskip}{5pt}\setlength{\parindent}{0pt}}

\fi

% ========================================================================

% Lists

\newcounter{ctr}
\newcounter{savectr}
\newcounter{ectr}
\newcounter{eectr}

\newenvironment{newenum}{%
\begin{list}{{\rm (\arabic{ctr})}\hfill}{\usecounter{ctr} \labelwidth=15pt%
\labelsep=7pt \leftmargin=22pt \topsep=0pt%
\setlength{\listparindent}{\saveparindent}%
\setlength{\parsep}{\saveparskip}%
\setlength{\itemsep}{2pt} }}{\end{list}}

\newenvironment{newenumbf}{%
\begin{list}{{\bf \arabic{ctr}.}\hfill}{\usecounter{ctr} \labelwidth=15pt%
\labelsep=5pt \leftmargin=20pt \topsep=2pt%
\setlength{\listparindent}{\saveparindent}%
\setlength{\parsep}{\saveparskip}%
\setlength{\itemsep}{2pt} }}{\end{list}}

\newenvironment{bignewenum}{%
\begin{list}{{\bf (\arabic{ctr})}\hfill}{\usecounter{ctr} \labelwidth=17pt%
\labelsep=7pt \leftmargin=24pt \topsep=4pt%
\setlength{\listparindent}{\saveparindent}%
\setlength{\parsep}{\saveparskip}%
\setlength{\itemsep}{4pt} }}{\end{list}}

\newenvironment{tiret}{%
\begin{list}{\hspace{0pt}\rule[0.5ex]{6pt}{1pt}\hfill}{\labelwidth=14pt%
\labelsep=1pt \leftmargin15pt \topsep=2pt%
\setlength{\listparindent}{\saveparindent}%
\setlength{\parsep}{\saveparskip}%
\setlength{\itemsep}{1pt} }}{\end{list}}

\newenvironment{onelist}{%
\begin{list}{{\rm (\arabic{ctr})}\hfill}{\usecounter{ctr} \labelwidth=20pt%
\labelsep=7pt \leftmargin=27pt \topsep=0pt%
\setlength{\listparindent}{\saveparindent}%
\setlength{\parsep}{\saveparskip}%
\setlength{\itemsep}{2pt}}}{\end{list}}

\newenvironment{twolist}{%
\begin{list}{{\rm (\arabic{ctr}.\arabic{ectr})}%
\hfill}{\usecounter{ectr} \labelwidth=24pt%
\labelsep=7pt \leftmargin=31pt \topsep=0pt%
\setlength{\listparindent}{\saveparindent}%
\setlength{\parsep}{\saveparskip}%
\setlength{\itemsep}{2pt} }}{\end{list}}

\newenvironment{newitemize}{%
\begin{list}{$\bullet$\hfill}{\labelwidth=10pt%
\labelsep=5pt \leftmargin=15pt \topsep=2pt%
\setlength{\listparindent}{\saveparindent}%
\setlength{\parsep}{\saveparskip}%
\setlength{\itemsep}{2pt} }}{\end{list}}

\newenvironment{steplist}{%
\begin{list}{\textbf{Step~{\arabic{ctr}:}}\hfill}
{\usecounter{ctr} \labelwidth=40pt%
\labelsep=5pt \leftmargin=45pt \topsep=5pt%
\setlength{\listparindent}{\saveparindent}%
\setlength{\parsep}{\saveparskip}%
\setlength{\itemsep}{3pt} }}{\end{list}}

\newenvironment{substeplist}{%
\begin{list}{\textbf{\arabic{ctr}.\arabic{ectr}}\hfill}
{\usecounter{ectr} \labelwidth=18pt%
\labelsep=5pt \leftmargin=23pt \topsep=2pt%
\setlength{\listparindent}{\saveparindent}%
\setlength{\parsep}{\saveparskip}%
\setlength{\itemsep}{2pt} }}{\end{list}}

\newlength{\savejot}
\setlength{\jot}{6pt}
\setlength{\savejot}{\jot}

\ifnum\llncs=0

\newenvironment{newmath}{\begin{displaymath}%
\setlength{\abovedisplayskip}{4pt}%
\setlength{\belowdisplayskip}{4pt}%
\setlength{\abovedisplayshortskip}{5pt}%
\setlength{\belowdisplayshortskip}{5pt} }{\end{displaymath}}

\newenvironment{neweqnarrays}{\begin{eqnarray*}%
\setlength{\abovedisplayskip}{-4pt}%
\setlength{\belowdisplayskip}{-4pt}%
\setlength{\abovedisplayshortskip}{-4pt}%
\setlength{\belowdisplayshortskip}{-4pt}%
\setlength{\jot}{-0.4in} }{\end{eqnarray*}}

\newenvironment{newequation}{\begin{equation}%
\setlength{\abovedisplayskip}{4pt}%
\setlength{\belowdisplayskip}{4pt}%
\setlength{\abovedisplayshortskip}{5pt}%
\setlength{\belowdisplayshortskip}{5pt} }{\end{equation}}

\else

\newenvironment{newmath}{\begin{displaymath}%
}{\end{displaymath}}
\newenvironment{neweqnarrays}{\begin{eqnarray*}%
}{\end{eqnarray*}}
\newenvironment{newequation}{\begin{equation}%
}{\end{equation}}

\fi

% =========================================================================

\ifnum\llncs=0

\makeatletter
\def\appearsin#1{\gdef\@appearsin{#1}}
\def\maketitle{\par
 \begingroup
 \def\thefootnote{\arabic{footnote}}
 % \def\thefootnote{\fnsymbol{footnote}}
 \def\@makefnmark{\hbox
 to 0pt{$^{\@thefnmark}$\hss}}
 \if@twocolumn
 \twocolumn[\@maketitle]
 \else \newpage
 \global\@topnum\z@ \@maketitle \fi\pagestyle{}\@thanks
 \endgroup
 \setcounter{footnote}{0}
 \let\maketitle\relax
 \let\@maketitle\relax
 \gdef\@thanks{}\gdef\@author{}\gdef\@title{}\gdef\@appearsin{}
          \let\thanks\relax}
\def\@maketitle{\newpage
 \noindent \@appearsin
 \vskip 0.5in \begin{center}
 {\LARGE \@title \par} \vskip 1.5em {\large \lineskip .5em
\begin{tabular}[t]{c}\@author
 \end{tabular}\par}
 \vskip 1em {\normalsize \@date} \end{center}
 \par
 \vskip 1.5em}


%\def\abstract{\if@twocolumn
%\section*{Abstract}
%\else \small
%\begin{center}
%{\bf Abstract\vspace{-.5em}\vspace{0pt}}
%\end{center}
%\quotation
%\fi}
%\def\endabstract{\if@twocolumn\else\endquotation\fi}
%\mark{{}{}}

\fi

% =========================================================================

% General

\ifnum\llncs=1
\newcommand{\mysubsubsection}[1]{\vspace{-5pt}\subsubsection{#1.}}
\newcommand{\headingb}[1]{\vspace{-5pt}\subsubsection{#1.}}
\else
\newcommand{\mysubsubsection}[1]{\subsubsection{#1}}
% \newcommand{\headingb}[1]{\subsection{#1}}
\newcommand{\headingb}[1]{{\vspace{5pt}\noindent\textbf{#1.}}}
\fi
\newcommand{\headingg}[1]{{\sc{#1}}}
\newcommand{\heading}[1]{{\vspace{5pt}\noindent\sc{#1}}}
%\newcommand{\vechint}{\mathbf{hint}}
%\newcommand{\hint}{\mathsf{hint}}

\newcommand{\nbbsie}{{\notionfont{nbbsie}}}


%\newcommand{\coin}{\mathsf{coin}}


\newcommand{\deffont}{\rm}
\newcommand{\schemefont}[1]{{\mathsf{#1}}}
\newcommand{\primfont}[1]{{\mathsf{#1}}}
\newcommand{\algfont}[1]{{\mathsf{#1}}}
\newcommand{\advfont}[1]{{\mathsf{#1}}}
\newcommand{\orfont}[1]{\mathsc{#1}}
\newcommand{\varfont}[1]{\mathit{#1}}
\newcommand{\randvarfont}[1]{\mathbf{#1}}
\newcommand{\constfont}[1]{\mathtt{#1}}
\newcommand{\notionfont}[1]{\mathrm{#1}}
\newcommand{\eventfont}[1]{{\mathsc{#1}}}
\newcommand{\vectorfont}[1]{\mathslbf{#1}}
\newcommand{\transformfont}[1]{\mathsf{#1}}

\newcommand{\PRG}{\schemefont{PRG}}

\newcommand{\iO}{\algfont{iO}}
\newcommand{\diO}{\algfont{diO}}
\newcommand{\AIPO}{\algfont{AIPO}}
\newcommand{\MBAIPO}{\algfont{MB\mbox{-}AIPO}}

\newcommand{\hashH}{\algfont{H}}
\newcommand{\hashG}{\algfont{G}}
%\newcommand{\PRF}{\algfont{F}}


\newcommand{\hyb}{\algfont{hy}}
\newcommand{\asy}{\algfont{asy}}
\newcommand{\sy}{\algfont{sy}}


\newcommand{\authnote}[2]{\ifnum\authnotes=1 \begin{center}\fbox{\begin{minipage}{5.7in}
\textbf{#1 says:} #2\end{minipage}}\end{center} \fi}

\newcommand{\main}{\procfont{Game}}


\ifnum\llncs=1
\newcommand{\gamesfontsize}{\footnotesize}
\else
\newcommand{\gamesfontsize}{\small}
\fi

% \footnotesize}
\ifnum\llncs=0
\renewcommand{\captionfont}{}
\else
%\newcommand{\captionfont}{}
\fi
\newcommand{\indd}{\hspace*{12pt}}

\newcommand{\mpage}[2]{\begin{minipage}{#1\textwidth} #2 \end{minipage}}
\newcommand{\fpage}[2]{\framebox{\begin{minipage}[t]{#1\textwidth}\setstretch{1.2}\gamesfontsize  #2 \end{minipage}}}

\newcommand{\hfpages}[3]{\hfpagess{#1}{#1}{#2}{#3}}
\newcommand{\hfpagess}[4]{
        \begin{tabular}{c@{\hspace*{.5em}}c}
        \framebox{\begin{minipage}[t]{#1\textwidth}\setstretch{1.2}\gamesfontsize #3 \end{minipage}}
        &
        \framebox{\begin{minipage}[t]{#2\textwidth}\setstretch{1.2}\gamesfontsize #4 \end{minipage}}
        \end{tabular}
    }
\newcommand{\hfpagesss}[6]{
        \begin{tabular}{c@{\hspace*{.1em}}c@{\hspace*{.1em}}c}
        \framebox{\begin{minipage}[t]{#1\textwidth}\setstretch{1.2}\gamesfontsize #4 \end{minipage}}
        &
        \framebox{\begin{minipage}[t]{#2\textwidth}\setstretch{1.2}\gamesfontsize #5 \end{minipage}}
        &
        \framebox{\begin{minipage}[t]{#3\textwidth}\setstretch{1.2}\gamesfontsize #6 \end{minipage}}
        \end{tabular}
    }

\newcommand{\hfpagessss}[8]{
        \begin{tabular}{c@{\hspace*{.5em}}c}
        \framebox{\begin{minipage}[t]{#1\textwidth}\setstretch{1.2}\gamesfontsize #5 \end{minipage}}
        &
        \framebox{\begin{minipage}[t]{#2\textwidth}\setstretch{1.2}\gamesfontsize
#6 \end{minipage}}\\[3in]
        \framebox{\begin{minipage}[t]{#3\textwidth}\setstretch{1.2}\gamesfontsize #7 \end{minipage}} 
        & 
		\framebox{\begin{minipage}[t]{#4\textwidth}\setstretch{1.2}\gamesfontsize #8 \end{minipage}}
        \end{tabular}
    }

\newcommand{\boxFourColTable}[8]{
		\fpage{0.98}{\smallskip
        \begin{tabular}{c@{\hspace*{.1em}}|@{\hspace*{.4em}}c@{\hspace*{.1em}}|@{\hspace*{.4em}}c@{\hspace*{.1em}}|@{\hspace*{.4em}}c@{\hspace*{.1em}}}
        {\begin{minipage}[t]{#1\textwidth}\setstretch{1}\gamesfontsize #5 \end{minipage}}
        &
        {\begin{minipage}[t]{#2\textwidth}\setstretch{1}\gamesfontsize #6 \end{minipage}}
        &
        {\begin{minipage}[t]{#3\textwidth}\setstretch{1}\gamesfontsize #7 \end{minipage}}
		&
		{\begin{minipage}[t]{#4\textwidth}\setstretch{1}\gamesfontsize #8 \end{minipage}}

        \end{tabular}\smallskip
        }
    }



\newcommand{\boxFourColThreeSepTable}[8]{
		\smallskip
        \begin{tabular}{c@{\hspace*{.1em}}|@{\hspace*{.4em}}c@{\hspace*{.1em}}|@{\hspace*{.4em}}c@{\hspace*{.2em}}c}
        {\begin{minipage}[t]{#1\textwidth}\setstretch{1}\gamesfontsize #5 \end{minipage}}
        &
        {\begin{minipage}[t]{#2\textwidth}\setstretch{1}\gamesfontsize #6 \end{minipage}}
        &
        {\begin{minipage}[t]{#3\textwidth}\setstretch{1}\gamesfontsize #7 \end{minipage}}
		&
		{\begin{minipage}[t]{#4\textwidth}\setstretch{1}\gamesfontsize #8 \end{minipage}}

        \end{tabular}\smallskip

    }

\newcommand{\hpagesl}[3]{\begin{tabular}{c|c}\begin{minipage}{#1\textwidth} #2 \end{minipage} & \begin{minipage}{#1\textwidth} #3 \end{minipage}\end{tabular}}
\newcommand{\hpagessl}[4]{\begin{tabular}{c|c}\begin{minipage}[t]{#1\textwidth} #3 \end{minipage} & \begin{minipage}[t]{#2\textwidth} #4 \end{minipage}\end{tabular}}
\newcommand{\hpages}[3]{\begin{tabular}{cc}\begin{minipage}[t]{#1\textwidth} #2 \end{minipage} & \begin{minipage}[t]{#1\textwidth} #3 \end{minipage}\end{tabular}}
\newcommand{\hpagess}[4]{
    \begin{tabular}{cc}
    \begin{minipage}[t]{#1\textwidth}\setstretch{1.1}\gamesfontsize  #3 \end{minipage}
    &
    \begin{minipage}[t]{#2\textwidth}\setstretch{1.1}\gamesfontsize #4 \end{minipage}\end{tabular}
   }
\newcommand{\hpagesss}[6]{
    \begin{tabular}{ccc}
    \begin{minipage}[t]{#1\textwidth}\setstretch{1.1}\gamesfontsize #4 \end{minipage} &
    \begin{minipage}[t]{#2\textwidth}\setstretch{1.1}\gamesfontsize #5 \end{minipage} &
    \begin{minipage}[t]{#3\textwidth}\setstretch{1.1}\gamesfontsize #6 \end{minipage}
    \end{tabular}}
\newcommand{\hpagessss}[8]{
    \begin{tabular}{cc}
    \begin{minipage}[t]{#1\textwidth}\setstretch{1.1}\gamesfontsize  #5 \end{minipage} &
    \begin{minipage}[t]{#2\textwidth}\setstretch{1.1}\gamesfontsize #6 \end{minipage}\\
    \begin{minipage}[t]{#3\textwidth}\setstretch{1.1}\gamesfontsize  #7 \end{minipage} & 
    \begin{minipage}[t]{#4\textwidth}\setstretch{1.1}\gamesfontsize  #8 \end{minipage}
    \end{tabular}}



\newcommand{\hpagesssl}[6]{
    \begin{tabular}{c|c|c}
    \begin{minipage}[t]{#1\textwidth} #4 \end{minipage} &
    \begin{minipage}[t]{#2\textwidth} #5 \end{minipage} &
    \begin{minipage}[t]{#3\textwidth} #6 \end{minipage}
    \end{tabular}}

\newcommand{\threeCols}[6]{
\begin{center}
        \framebox{
        \begin{tabular}{@{\hspace{-0.2em}}c@{\hspace{0.2em}}|@{\hspace{0.2em}}c@{\hspace{0.2em}}|@{\hspace{0.2em}}c@{\hspace{0.2em}}}
        \begin{minipage}[t]{#1\textwidth}\setstretch{1.1}\gamesfontsize #4
        \end{minipage} &
        \begin{minipage}[t]{#2\textwidth}\setstretch{1.1}\gamesfontsize #5
        \end{minipage} &
        \begin{minipage}[t]{#3\textwidth}\setstretch{1.1}\gamesfontsize #6
        \end{minipage}
        \end{tabular}
        }
\end{center}
}

\newcommand{\twoCols}[4]{
\begin{center}
        \framebox{
        \begin{tabular}{c@{\hspace*{.4em}}|c@{\hspace*{.4em}}c}
        \begin{minipage}[t]{#1\textwidth}\setstretch{1.2}\gamesfontsize #3 \end{minipage}
        &
        \begin{minipage}[t]{#2\textwidth}\setstretch{1.2}\gamesfontsize #4 \end{minipage}
        \end{tabular}
        }
\end{center}
}


\newcommand{\gameName}[1]{\underline{$\main$ #1} \\[2pt]}
\newcommand{\procName}[1]{\underline{#1} \\[2pt]}

\newcommand{\bnm}{\begin{newmath}}
\newcommand{\enm}{\end{newmath}}

\newcommand{\bea}{\begin{eqnarray*}}
\newcommand{\eea}{\end{eqnarray*}}


\newcommand{\Pub}{\mathsf{P}}
\newcommand{\Priv}{\mathsf{p}}
\newcommand{\FF}{\mathsf{FF}}
\newcommand{\tuple}[1]{\langle{#1}\rangle}
\def\from{\mbox{from}\ }
\def\From{\mbox{From}\ }
\def\bits{\{0,1\}}
\def\cross{\times}
\newcommand{\xor}{\,{\oplus}\,}
\newcommand{\Colon}{{:\;\;}}
\newcommand{\mystrut}{\rule{0em}{15pt}}
\newcommand{\namestrut}{\rule{0em}{20pt}}
\def\poly{\mathop{\rm poly}\nolimits}
\def\emptystring{\varepsilon}
\def\emptyid{()}
\newcommand{\Dom}{\mathsf{Dom}}
\newcommand{\DomR}{\mathsf{DomR}}
\newcommand{\Rng}{\mathsf{Rng}}
\newcommand{\RngR}{\mathsf{RngR}}
\newcommand{\calA}{{\cal A}}
\newcommand{\calB}{{\cal B}}
\newcommand{\calC}{{\cal C}}
\newcommand{\calF}{{\cal F}}
\newcommand{\calK}{{\cal K}}
\newcommand{\calM}{{\cal M}}
\newcommand{\calO}{{\cal O}}
\newcommand{\calR}{{\cal R}}
\newcommand{\calH}{{\cal H}}
\newcommand{\calG}{{\cal G}}
\newcommand{\calD}{{\cal D}}
\newcommand{\calE}{{\cal E}}
\newcommand{\calP}{{\cal P}}
\newcommand{\calS}{{\cal S}}
\newcommand{\calT}{{\cal T}}
\newcommand{\calU}{{\cal U}}
\newcommand{\calX}{{\cal X}}
\newcommand{\calW}{{\cal W}}
\newcommand{\calY}{{\cal Y}}
\newcommand{\calV}{{\cal V}}
\newcommand{\calEE}{{\cal EE}}
\newcommand{\N}{{{\mathbb N}}}
\newcommand{\Z}{{{\mathbb Z}}}
\newcommand{\R}{{{\mathbb R}}}
\newcommand{\goesto}{{\rightarrow}}
\newcommand{\eqdef}{\;\stackrel{\rm def}{=}\;}
\newcommand{\then}{{\;;\;\;}}   % for [ ; ; ; : ] notation
\newcommand{\andthen}{{\;:\;\;}}
\def\union{\cup}
\def\bigunion{\bigcup}
\def\intersection{\cap}
\def\bigintersection{\bigcap}
\def\suchthatt{\: :\:}
\newcommand{\suchthat}{{\mbox{s.t.\ }}}
\def\next{\,;\,}
\newcommand{\sett}[1]{\{#1\}}
\newcommand{\set}[2]{\{#1 \suchthatt #2\}}
\newcommand{\setsize}[1]{\left|{#1}\right|}
\def\leqq{\;\leq\;}
\def\eqq{\;=\;}
\def\geqq{\;\geq\;}
\def\equivv{\;\equiv\;}
\def\prn#1{\left(#1\right)}
% \def\getsr{\stackrel{{\scriptscriptstyle R}}{\leftarrow}}
\newcommand{\getsr}{{{\,\leftarrow{\hspace*{-3pt}\raisebox{.75pt}{$\scriptscriptstyle\$$}}\,}}}
% \renewcommand{\choose}[2]{{{#1}\atopwithdelims(){#2}}}
\newcommand{\OR}{{\textstyle\bigvee}}
\newcommand{\Var}{{\mbox{\bf Var}}}
\newcommand{\E}{\mathbf{E}}
\newcommand{\EE}[1]{{\E\left[{#1}\right]}}
\newcommand{\EEE}[2]{{\E_{#1}\left[{#2}\right]}}
\newcommand{\Prob}[1]{{\Pr\left[\,{#1}\,\right]}}
\newcommand{\probb}[2]{{\Pr}_{#1}\left[\,{#2}\,\right]}
\newcommand{\probbs}[3]{{\Pr}^{#1}_{#2}\left[\,{#3}\,\right]}
\newcommand{\probbp}[2]{{\Pr}_{#1}'\left[\,{#2}\,\right]}
\newcommand{\Probb}[2]{{{\Pr}_{#1}\left[\,{#2}\,\right]}}
\newcommand{\Probbp}[2]{{{\Pr}_{#1}'\left[\,{#2}\,\right]}}
\newcommand{\condProb}[2]{{\Pr}\left[\,#1\,|\,#2\,\right]}
\newcommand{\CondProb}[2]{{\Pr}\left[\: #1\:\left|\right.\:#2\:\right]}
\newcommand{\CondProbb}[3]{{\Pr}_{#1}\left[\: #2\:\left|\right.\:#3\:\right]}
\newcommand{\CondProbbp}[3]{%
{\Pr}_{#1}^{'}\left[\: #2\:\left|\right.\:#3\:\right]}
\ifnum\fullversion=1
\newcommand{\ProbExp}[2]{{\Pr}\left[\: #1\:\suchthatt\:#2\:\right]}
\else
\newcommand{\ProbExp}[2]{{\Pr}[#1\,:\,#2]}
\fi
\newcommand{\expect}[1]{{{\mathbf{E}}[{#1}]}}
\newcommand{\expectt}[1]{{{\mathbf{E}}\left[{#1}\right]}}
\newcommand{\qquadd}{{\quad}}
\def\d{{\delta}}
\def\e{{\epsilon}}
\newcommand{\sfrac}[2]{{\textstyle\frac{#1}{#2}}}
\newcommand{\ceiling}[1]{\lceil #1\rceil}
\def\pprod{{\textstyle\prod}}
\def\ssum{{\textstyle\sum}}
\newcommand{\concat}{{\,\|\,}}
\newcommand{\infi}[2]{\inf_{#1}\left\{#2\right\}}
\newcommand{\supr}[2]{\sup_{#1}\left\{#2\right\}}
\newcommand{\Maps}{\mathsf{Maps}}

\newcommand{\dcomp}{\mathsf{dist}\mbox{-}\mathsf{comp}}

% ========================================================================

\newcommand{\verylongleftarrow}[1]
      {\setlength{\unitlength}{.01in}
      \begin{picture}(#1,1) \put(#1,0){\vector(-1,0){#1}} \end{picture}}
\newcommand{\verylongrightarrow}[1]             %longleft and rightgoing arrows
      {\setlength{\unitlength}{.01in}           %for protocols
      \begin{picture}(#1,1) \put(0,0){\vector(1,0){#1}} \end{picture}}
\newcommand{\leftgoing}[2]{{\stackrel{{\displaystyle #2}} {\verylongleftarrow{#1}}}}
\newcommand{\rightgoing}[2]{{\stackrel{{\displaystyle #2}} {\verylongrightarrow{#1}}}}

\newcommand{\leftgoinga}[1]{\leftgoing{230}{#1} }
\newcommand{\rightgoinga}[1]{\rightgoing{230}{#1} }

\newcommand{\leftgoingb}[1]{\leftgoing{300}{#1} }
\newcommand{\rightgoingb}[1]{\rightgoing{300}{#1} }

\newcommand{\seq}{\!=\!}
\newcommand{\sminus}{\!-\!}
\newcommand{\splus}{\!+\!}
\def\next{\:;\:}

% ===================================================================
\newcommand{\kwfont}[1]{{\ensuremath{\mathrm{#1}}}}
\newcommand{\kwfunction}{{\kwfont{function}\ }}
\newcommand{\kwlabel}{{\kwfont{label}\ }}
\newcommand{\kwfor}{{\kwfont{for}\ }}
\newcommand{\kwand}{{\kwfont{and}\ }}
\newcommand{\kwor}{{\kwfont{or}\ }}
\newcommand{\kwnot}{{\kwfont{not}}}
\newcommand{\kwdo}{{\kwfont{do}\ }}
\newcommand{\kwreturn}{{\kwfont{return}\ }}
\newcommand{\kwReturn}{{\kwfont{Return}\ }}
\newcommand{\kwalgorithm}{{\ensuremath{\mathbf{Algorithm}\ }}}
\newcommand{\kwprotocol}{{\kwfont{Protocol}\ }}
\newcommand{\kwexperiment}{{\kwfont{Experiment}\ }}
\newcommand{\kwadversary}{{\kwfont{Adversary}\ }}
\newcommand{\kworacle}{{\kwfont{Oracle}\ }}
\newcommand{\kwuntil}{{\kwfont{until}\ }}
\newcommand{\kwrepeat}{{\kwfont{repeat}\ }}
\newcommand{\kwif}{{\kwfont{If}\ }}
\newcommand{\kwthen}{{\kwfont{Then}\ }}
\newcommand{\kwelse}{{\kwfont{Else}\ }}
\newcommand{\kwabort}{{\kwfont{abort}\ }}
\newcommand{\kwgoto}{{\kwfont{goto}\ }}
\newcommand{\kwwhile}{{\kwfont{while}\ }}
\newcommand{\kwparse}{{\kwfont{parse}\ }}
\newcommand{\kwas}{{\kwfont{as}\ }}
\newcommand{\kwstatic}{{\kwfont{static}\ }}
\newcommand{\kwrun}{{\kwfont{run}\ }}
\newcommand{\kwbegin}{{\kwfont{begin}\ }}
\newcommand{\kwend}{{\kwfont{end}\ }}
\newcommand{\kwstart}{{\kwfont{start}\ }}
\newcommand{\kwcontinue}{{\kwfont{continue}\ }}
\newcommand{\kwdefine}{{\kwfont{define}\ }}
\newcommand{\kwflip}{{\kwfont{flip}\ }}
\newcommand{\kwlet}{{\kwfont{let}\ }}
\newcommand{\kwof}{{\kwfont{of}\ }}
\newcommand{\kwcase}{{\kwfont{case}\ }}
\newcommand{\kwswitch}{{\kwfont{switch}\ }}
\newcommand{\kwpick}{{\kwfont{pick}\ }}
\newcommand{\kwset}{{\kwfont{set}\ }}
\newcommand{\kwcompute}{{\kwfont{compute}\ }}
\newcommand{\comment}[1]{\hspace{15pt}{\small /$\!\!$/\ #1}}
\newcommand{\Comment}[1]{\hspace{5pt}{/$\!\!$/\ #1}}


% ==================================================================
% Definitions for this paper
% ==================================================================

\newcommand{\Coins}{\mathsl{Coins}}
\newcommand{\msgsp}{\mathsl{MsgSp}}
\newcommand{\tagsp}{\mathsl{TagSp}}

\newcommand{\XGame}[3]{\mathrm{#1}_{#2}^{#3}(\lambda)}
\newcommand{\vk}{\mathit{vk}}
\newcommand{\pars}{\params}
\newcommand{\ek}{\mathsl{ek}}
\newcommand{\td}{\mathsl{td}}
\newcommand{\pk}{\mathsl{pk}}
\newcommand{\pp}{\mathsl{pp}}
\newcommand{\sk}{\mathsl{sk}}
\newcommand{\dk}{\mathsl{dk}}
\newcommand{\ak}{\mathsl{ak}}
\newcommand{\hk}{\mathsl{hk}}
%\renewcommand{\AE}{\Pi}
%\newcommand{\enc}{\mathcal{E}}
\newcommand{\keys}{\mathcal{K}}
\newcommand{\keysh}{\agen}
%\newcommand{\decc}{\mathcal{D}}
\newcommand{\pgen}{\mathcal{P}}
\newcommand{\agen}{\mathcal{A}}
\newcommand{\verify}{\mathcal{V}}
\newcommand{\dverify}{\mathcal{W}}
\newcommand{\Hash}{H}
\newcommand{\aux}{a}
\newcommand{\hf}{\Gamma}

%\newcommand{\schemefont}[1]{{\mathsf{#1}}}
%\newcommand{\algfont}[1]{{\mathsf{#1}}}
\newcommand{\Adv}{\mathbf{Adv}}
\newcommand{\AdvGen}[3]{\mathbf{Adv}^{\mathrm{#1}}_{#2}(#3)}
\newcommand{\Exp}{\mathbf{Exp}}
\newcommand{\Op}{\mathbf{Op}}
\newcommand{\advclass}[2]{\mathcal{A}^{#1}_{#2}}
\newcommand{\advA}{A}
\newcommand{\advB}{B}
\newcommand{\advI}{I}
\newcommand{\advW}{W}
\newcommand{\advH}{H}
\newcommand{\advU}{U}
\newcommand{\advX}{C}
\newcommand{\advS}{S}
\newcommand{\advE}{E}
\newcommand{\INDCPA}{\mbox{IND-CPA}}
\newcommand{\ExpINDCPA}[2]{\mathrm{INDCPA}_{#1,#2}}
\newcommand{\AdvINDCPA}[2]{\AdvGen{indcpa}{#1}{#2}}
\newcommand{\AdvSOAR}[2]{\mathbf{Adv}^{\mathrm{soa}\mbox{-}\mathrm{c}}_{#1}(#2)}
\newcommand{\AdvSOAK}[2]{\mathbf{Adv}^{\mathrm{soa}\mbox{-}\mathrm{k}}_{#1}(#2)}


\newcommand{\atk}{\mathrm{atk}}
\newcommand{\cpa}{\mathrm{cpa}}
\newcommand{\cca}{\mathrm{cca}}
\newcommand{\ind}{\mathrm{ind}}

\newcommand{\rv}[1]{\mathsf{#1}}

\newcommand{\MsgSp}[1]{\msgspace_{#1}}
\newcommand{\Comp}[1]{\overline{#1}}
\newcommand{\BADc}{\eventfont{Bad}}

\addtolength{\belowcaptionskip}{-2mm}
\addtolength{\abovecaptionskip}{-3mm}
\addtolength{\topsep}{-1mm}

\newcommand{\procfont}[1]{\mathsc{#1}}
\newcommand{\tablefont}[1]{\mathsf{#1}}
\newcommand{\Initialize}{\procfont{Initialize}}
\newcommand{\Finalize}{\procfont{Finalize}}
\newcommand{\procedure}[2]{\underline{\procfont{proc}\ #1(#2)}}
\newcommand{\Onquery}[2]{\underline{\procfont{#1}(#2)}}
\newcommand{\InitializeCode}[1]{\Onquery{Initialize}{#1}}
\newcommand{\FinalizeCode}[1]{\Onquery{Finalize}{#1}}

\newcommand{\EncO}{\procfont{Enc}}
\newcommand{\DecO}{\procfont{Dec}}
\newcommand{\MsgO}{\procfont{Msg}}
\newcommand{\KDO}{\procfont{Kd}}
\newcommand{\OpO}{\procfont{Op}}
%\newcommand{\FO}{\procfont{F}}
\newcommand{\LR}{\procfont{LR}}
\newcommand{\CorruptO}{\procfont{Corrupt}}

\newcommand{\dsim}{\overline{S}}
\newcommand{\st}{\mathit{St}}
\newcommand{\newst}{\overline{\mathit{St}}}
\newcommand{\group}{\mathbb{G}}
\newcommand{\groupp}{\mathbb{H}}
\newcommand{\goutputs}{\:{\Rightarrow}\:}
\newcommand{\outputs}{\textrm{ outputs }}
\newcommand{\wins}{\textrm{ wins}}
\newcommand{\true}{\mathsf{true}}
\newcommand{\false}{\mathsf{false}}
\newcommand{\AND}{\;\wedge\;}
\newcommand{\bad}{\mathsf{bad}}
\newcommand{\Good}{\mathsf{Good}}
\newcommand{\Gm}{\textnormal{G}}
\newcommand{\fn}{\footnotesize}

\newcommand{\params}{\pi}

\newcommand{\CDL}{\mathsf{CDL}}

\newcommand{\Sparkle}{\mathsf{Sparkle}+}

\newcommand{\vecb}{\mathbf{b}}
\newcommand{\vecg}{\mathbf{g}}
\newcommand{\vech}{\mathbf{h}}
\newcommand{\veca}{\mathbf{a}}
\newcommand{\vecc}{\mathbf{c}}
\newcommand{\vecd}{\mathbf{d}}
\newcommand{\vece}{\mathbf{e}}
\newcommand{\veci}{\mathbf{i}}
\newcommand{\vecm}{\mathbf{m}}
\newcommand{\vecell}{\mathbf{\ell}}
\newcommand{\vecp}{\mathbf{p}}
\newcommand{\veco}{\mathbf{o}}
\newcommand{\vecq}{\mathbf{q}}
\newcommand{\vecr}{\mathbf{r}}
\newcommand{\vecro}{\overline{\mathbf{r}}}
\newcommand{\vecs}{\mathbf{s}}
\newcommand{\vecv}{\mathbf{v}}
\newcommand{\vecK}{\mathbf{K}}
\newcommand{\vecw}{\mathbf{w}}
\newcommand{\vecmo}{\overline{\mathbf{m}}}
\newcommand{\vecms}{\vecm^*}
\newcommand{\vecwo}{\overline{\mathbf{w}}}
\newcommand{\vecsko}{\overline{\mathbf{dk}}}
\newcommand{\vecy}{\mathbf{y}}
\newcommand{\vecz}{\mathbf{z}}
\newcommand{\vecx}{\mathbf{x}}
\newcommand{\vecxo}{\overline{\mathbf{x}}}
\newcommand{\vecxs}{\vecx^*}
\newcommand{\vecbot}{\mathbf{\bot}}
\newcommand{\vecpk}{\mathbf{pk}}
\newcommand{\vecsk}{\mathbf{sk}}
\newcommand{\vecu}{\mathbf{u}}
\newcommand{\vect}{\mathbf{t}}
\newcommand{\vecalpha}{{\boldsymbol\alpha}}
\newcommand{\vecbeta}{{\boldsymbol\beta}}
\newcommand{\matM}{\mathbf{A}}
\newcommand{\vecA}{\mathbf{A}}

\newcommand{\NR}{r}
\newcommand{\NS}{i}
\newcommand{\SOAout}{w}
\newcommand{\subAlg}{F}
\newcommand{\state}{St}
\newcommand{\PrMC}{U}



\newcommand{\G}{\mathbb{G}}
\newcommand{\Zp}{\mathbb{Z}_p}
\renewcommand{\P}{\mathcal{P}}
\newcommand{\V}{\mathcal{V}}

\newcommand{\vecpkI}[1]{\mathbf{ek}[#1]}
\newcommand{\vecskI}[1]{\mathbf{dk}[#1]}
\newcommand{\veccI}[1]{\mathbf{c}[#1]}
\newcommand{\vecmI}[1]{\mathbf{m}[#1]}
\newcommand{\vecrI}[1]{\mathbf{r}[#1]}
\newcommand{\recI}[1]{\mathbf{rec}[#1]}


\newcommand{\vecpkS}[1]{\mathbf{ek}_{#1}}
\newcommand{\vecmS}[1]{\mathbf{m}_{#1}}
\newcommand{\veccS}[1]{\mathbf{c}_{#1}}
\newcommand{\recS}[1]{\mathbf{rec}_{#1}}
\newcommand{\CRS}[1]{\mathit{CR}_{#1}}
\newcommand{\CSS}[1]{\mathit{CS}_{#1}}
\newcommand{\NRS}[1]{r_{#1}}
\newcommand{\NSS}[1]{i_{#1}}
\newcommand{\SOAoutS}[1]{w_{#1}}
\newcommand{\advSS}[1]{S_{#1}}


\newcommand{\vecpkSI}[2]{\mathbf{ek}_{#1}[#2]}
\newcommand{\vecskSI}[2]{\mathbf{dk}_{#1}[#2]}
\newcommand{\vecmSI}[2]{\mathbf{m}_{#1}[#2]}
\newcommand{\vecrSI}[2]{\mathbf{r}_{#1}[#2]}
\newcommand{\veccSI}[2]{\mathbf{c}_{#1}[#2]}
\newcommand{\recSI}[2]{\mathbf{rec}_{#1}[#2]}




\newcommand{\AdvCR}[2]{\mathbf{Adv}^{\mathrm{col}}_{#1}(#2)}
\newcommand{\rel}{\mathsf{R}}
\newcommand{\sam}{\mathcal{M}}
\newcommand{\ComEnc}{{\mathrm{com}}}
\newcommand{\AdvComEnc}[2]{\Adv^{\ComEnc}_{#1}(#2)}
\newcommand{\SEMSOA}{\ensuremath{\mathrm{SEM\mbox{-}SOA}}}
\newcommand{\SEMSOAREAL}{\ensuremath{\mathrm{SEM\mbox{-}SOA\mbox{-}REAL}}}
\newcommand{\SEMSOAIDEAL}{\ensuremath{\mathrm{SEM\mbox{-}SOA\mbox{-}IDEAL}}}
\newcommand{\AdvSEMSOA}[2]{\Adv^{\SEMSOA}_{#1}(#2)}
\newcommand{\SCR}{{\mathrm{scr}}}
\newcommand{\AdvSCR}[2]{\Adv^{\SCR}_{#1}(#2)}
\newcommand{\ExpSCR}[2]{\Exp^{\SCR}_{#1}(#2)}
\newcommand{\plaintext}{\mathrm{PT}}
\newcommand{\ciphertext}{\mathrm{CT}}
\newcommand{\CR}{\mathit{CR}}
\newcommand{\CS}{\mathit{CS}}
\newcommand{\ext}{\mathrm{Emb}}

\newcommand{\vlist}[2]{\langle #1:#2\rangle}
\newcommand{\NP}{\mathbf{NP}}
\newcommand{\accprob}[2]{\mathbf{AP}_{#1}(#2)}

\newcommand{\ksp}[2]{\mathsf{Keys}_{#1}(#2)}
\newcommand{\msp}[2]{\mathsf{Msgs}_{#1}(#2)}

% ==================================================================
% From Brent's Header
% ==================================================================
\newcommand{\G}{\ensuremath{\mathbb{G}}}
\newcommand{\Zp}{\ensuremath{{\Z_p}}}


% ==================================================================
% Stuff
% ==================================================================
\newcommand{\secpar}{\lambda}
\newcommand{\FEnc}{\mathsf{FE}}
\newcommand{\PEnc}{\mathsf{PE}}
\newcommand{\setup}{\mathsf{Setup}}
\newcommand{\kder}{\mathsf{KDer}}


\newcommand{\Add}{\mathsf{Add}}
\newcommand{\Val}{\mathsf{Val}}
\newcommand{\Setup}{\mathsf{Setup}}



\newcommand{\TBTDF}{\mathsf{TB\mbox{-}TDF}}
\newcommand{\NIZK}{\mathsf{NIZK}}
\newcommand{\TDF}{\mathsf{TDF}}
\newcommand{\PKE}{\mathsf{PKE}}
\newcommand{\TBPKE}{\mathsf{TB}\mbox{-}\mathsf{PKE}}
\newcommand{\kg}{\mathsf{Kg}}
\newcommand{\Kg}{\mathsf{Kg}}
\newcommand{\tg}{\mathsf{Tdg}}
\newcommand{\enc}{\mathsf{Enc}}
\newcommand{\decc}{\mathsf{Dec}}



\newcommand{\SKE}{\mathsf{SKE}}
\newcommand{\aee}{\mathsf{AE}}

\newcommand{\skg}{\mathsf{kg}}
\newcommand{\senc}{\mathsf{enc}}
\newcommand{\sdec}{\mathsf{dec}}


\newcommand{\KEM}{\mathsf{KEM}}
\newcommand{\kemkg}{\mathsf{KEM}.\mathsf{kg}}
\newcommand{\kemenc}{\mathsf{KEM}.\mathsf{enc}}
\newcommand{\kemdec}{\mathsf{KEM}.\mathsf{dec}}
\newcommand{\encaps}{\mathsf{Encaps}}
\newcommand{\decaps}{\mathsf{Decaps}}


\newcommand{\unq}{\mathsf{unq}}
\newcommand{\valid}{\mathsf{Valid}}
\newcommand{\MAC}{\mathsf{MAC}}
\newcommand{\tagg}{\mathsf{mac}}  %% GF: I changed this, to avoid confusion when MACs are used with tag-based primitives
\newcommand{\Tag}{\mathsf{Tag}}
\newcommand{\ver}{\mathsf{ver}}
\newcommand{\Ver}{\mathsf{Ver}}
\newcommand{\Prov}{\mathsf{Prov}}
\newcommand{\hc}{\mathsf{hc}}

\newcommand{\flag}{\mathsf{flag}}

\newcommand{\rec}{\mathsf{Rec}}

\newcommand{\atdf}{\mathsf{ATDF}}
\newcommand{\eval}{\mathsf{Eval}}
\newcommand{\inv}{\mathsf{Inv}}
\newcommand{\tbatdf}{\mathsf{TB}\mbox{-}\mathsf{ATDF}}

\newcommand{\dc}{\bot}
\newcommand{\parse}{\mathsf{Parse}}
\newcommand{\setsbad}{~\text{sets}~\mathsf{bad}}
\newcommand{\el}{\mathsf{El}}
\newcommand{\enct}{\mathsf{enc}}
\newcommand{\kdt}{\mathsf{kd}}
\newcommand{\msgt}{\mathsf{msg}}
\newcommand{\opt}{\mathsf{op}}
\newcommand{\Col}{\mathsf{Col}}
%\newcommand{\func}[1]{\calF_{\mathrm{#1}}}
\newcommand{\predi}[1]{\calP_{\mathrm{#1}}}
%\newcommand{\funcc}[2]{\calF_{\mathrm{#1}}^{#2}}
\newcommand{\id}{\mathsf{id}}
\newcommand{\msglen}{\mathsf{msglen}}
\newcommand{\eform}{\mathbf{E}}
\newcommand{\efunc}{\calF_\mathbf{e}}
\renewcommand{\epsilon}{\varepsilon}

\renewcommand{\prec}{\mathsf{pRec}}
\renewcommand{\cons}{\mathsf{Cons}}
\newcommand{\inval}{\mathsf{Inval}}

\newcommand{\obf}{\mathsf{Obf}}


\newcommand{\deThreeScheme}{\schemefont{DE3}}
\newcommand{\DEthreeKg}{\schalg{DE3}{Kg}}
\newcommand{\DEthreeEnc}{\schalg{DE3}{Enc}}
\newcommand{\DEthreeDec}{\schalg{DE3}{Dec}}
\newcommand{\func}{\delta}
\newcommand{\idfunc}{\mathsl{id}}
\newcommand{\xxx}{\mathrm{xxx}}
\newcommand{\neGuess}{\neScheme.\schemefont{Guess}}
\newcommand{\funcc}{\overline{\func}}


\newcommand{\padOAEP}{
\psset{arrowsize=0.2}
\begin{psmatrix}[colsep=0.5cm,rowsep=0.1cm]
% &\textrm{OAEP}\psspan{3}\\[2ex]
&m \in\{0,1\}^\ms & & r \in \{0,1\}^\rs\\[-1.5ex]
&[name=m]&&[name=r]&\\[1.0ex]
&[name=p1]\oplus&[name=H1]\psframebox{G}&[name=x1]\\[1.0ex]
&[name=x2]&[name=H2]\psframebox{H}&[name=p2]\oplus\\[1.0ex]
&[name=s]&&[name=t]\\[-1.5ex]
&s&&t
\end{psmatrix}
\ncline{->}{m}{s}
\ncline{->}{r}{t}
\ncline{-}{p1}{H1}
\ncline{-}{x1}{H1}
\ncline{-}{p2}{H2}
\ncline{-}{x2}{H2}
}





%%% Local Variables: 
%%% mode: latex
%%% TeX-master: "main-full"
%%% End: 

%\newcommand{\nbbsie}{{\notionfont{nbbsie}}}

%\newcommand{\iO}{\algfont{iO}}
%\newcommand{\diO}{\algfont{diO}}
%\newcommand{\AIPO}{\algfont{AIPO}}
%\newcommand{\MBAIPO}{\algfont{MB\mbox{-}AIPO}}

%\newcommand{\hashH}{\algfont{H}}
%\newcommand{\hashG}{\algfont{G}}
%\newcommand{\PRF}{\algfont{F}}


%\newcommand{\hyb}{\algfont{hy}}
%\newcommand{\asy}{\algfont{asy}}
%\newcommand{\sy}{\algfont{sy}}


%\newcommand{\FDom}{\mathsf{FDom}}
%\newcommand{\FRng}{\mathsf{FRng}}

\newcommand{\negl}{\schemefont{negl}}

\newcommand{\PRF}{\schemefont{PRF}}
\newcommand{\PRFK}{\schalg{\PRF}{Kg}}
\newcommand{\PRFE}{\schalg{\PRF}{Eval}}
\newcommand{\PRFPunct}{\schalg{\PRF}{Punct}}
\newcommand{\PRFDom}{\schalg{PRF}{Dom}}
\newcommand{\PRFRng}{\schalg{PRF}{Rng}}

\newcommand{\PRFGame}{\ensuremath{\textrm{PRF-DIST}}}
\newcommand{\pprf}{{\notionfont{pprf}}}

\newcommand{\AIPOGame}{\ensuremath{\textrm{AIPO}}}
\newcommand{\MBAIPOGame}{\ensuremath{\textrm{MB\mbox{-}AIPO}}}

\newcommand{\aipo}{{\notionfont{aipo}}}
\newcommand{\mbaipo}{{\notionfont{mb\mbox{-}aipo}}}
\newcommand{\dio}{{\notionfont{dio}}}

\newcommand{\lrow}{{\notionfont{lrowf}}}
\newcommand{\nm}{{\notionfont{nm}}}


\newcommand{\ELF}{\schemefont{ELF}}
\newcommand{\ELFikg}{\schalg{\ELF}{IKg}}
\newcommand{\ELFlkg}{\schalg{\ELF}{LKg}} 
\newcommand{\ELFeval}{\schalg{\ELF}{Eval}}
\newcommand{\ELFinv}{\schalg{\ELF}{Inv}} 
\newcommand{\ELFkl}{\schalg{\ELF}{kl}} 
\newcommand{\ELFDom}{\schalg{ELF}{Dom}}
\newcommand{\ELFRng}{\schalg{ELF}{Rng}}

\newcommand{\PRGK}{\schalg{\PRG}{Kg}}
\newcommand{\PRGE}{\schalg{\PRG}{Eval}}
\newcommand{\PRGDom}{\schalg{PRG}{Dom}}
\newcommand{\PRGRng}{\schalg{PRG}{Rng}}

\newcommand{\elf}{{\notionfont{elf}}}
\newcommand{\FODec}{\schalg{FO}{Dec}}
\newcommand{\FOKg}{\schalg{FO}{Kg}}
\newcommand{\FOEnc}{\schalg{FO}{Enc}}
\newcommand{\FO}{\schemefont{FO}}

\newcommand{\NFODec}{\schalg{\overline{FO}}{Dec}}
\newcommand{\NFOKg}{\schalg{\overline{FO}}{Kg}}
\newcommand{\NFOEnc}{\schalg{\overline{FO}}{Enc}}
\newcommand{\NFO}{\schemefont{\overline{FO}}}


\newcommand{\UFCPA}{\ensuremath{\notionfont{UF\mbox{-}CPA}}}

\newcommand{\INTCTXT}{\ensuremath{\notionfont{INT\mbox{-}CTXT}}}

%!TEX root = main.tex

%----EDITING----------------

%----MATH-----------------

% Define symbol
\newcommand{\defeq}{\coloneqq}

% Math hyphen
\mathchardef\mhyphencar ="2D
\newcommand{\dash}{\mhyphencar}%{\ensuremath{\mhyphencar\allowbreak}}

% Allow long equations to break across pages
\allowdisplaybreaks

% Punctuation in displaymath
\newcommand{\dmp}{\thinspace \text{.}}
\newcommand{\dmc}{\thinspace \text{,}}
\newcommand{\dmsc}{\thinspace \text{;}}

% Vectors in boldface
\renewcommand{\vec}[1]{\mathbf{#1}}
\newcommand{\gvec}[1]{\boldsymbol #1}

% Sets
\newcommand{\NN}{\mathbb{N}}
\newcommand{\ZZ}{\mathbb{Z}}
\newcommand{\RR}{\mathbb{R}}
\newcommand{\GG}{\mathbb{G}}
\newcommand{\SetS}{\mathsf{S}}
\newcommand{\SetT}{\mathsf{T}}

% Functions
\newcommand{\Func}{\mathrm{Fun}}
\newcommand{\Inj}{\mathrm{Inj}}
% Fractions
\newcommand{\num}[1]{\hat{#1}}
\newcommand{\den}[1]{\check{#1}}

% Least common multiple
\DeclareMathOperator{\lcm}{lcm}

% Variables appearing in a polynomial
\newcommand{\var}[1]{T_{#1}}


% Spaced dots
\newcommand{\ds}{\;\!}
\newcommand{\spaceddots}{. \ds . \ds .}


%----NUMBERING OF EQUATIONS-----------------

\newcommand{\Dlog}{\mathsf{Dlog}}
% Ture/false

% Sample
\newcommand{\sample}{\mathrel{\leftarrow\mkern-14mu\leftarrow}} 

% Bits and strings
\newcommand{\bin}{\{0,1\}}
\newcommand{\str}{\bin^*}
\newcommand{\secparam}{1^\secpar}
% Advantage

% Adversaries
\newcommand{\adv}{\mathcal{A}}
\newcommand{\advb}{\mathcal{B}}
\newcommand{\advc}{\mathcal{C}}
\newcommand{\advd}{\mathcal{D}}
\newcommand{\advs}{\mathcal{S}}
\newcommand{\advf}{\mathcal{F}}
\newcommand{\adve}{\mathcal{E}}
\newcommand{\advp}{\mathcal{P}}
\newcommand{\advq}{\mathcal{Q}}
\newcommand{\advr}{\mathcal{R}}
\newcommand{\advt}{\mathcal{T}}
\newcommand{\advu}{\mathcal{U}}
\newcommand{\advx}{\mathcal{X}}
\newcommand{\advz}{\mathcal{Z}}

% Bit-fixing adversaries
\newcommand{\algA}{\mathfrak{A}}
\newcommand{\algD}{\mathfrak{D}}

% I/O of 
% Various tables and lists
\newcommand{\List}{\mathsf{L}} 
\newcommand{\plist}{\mathcal{L}}
\newcommand{\tableS}{T_\encode}
\newcommand{\tableT}{T_\ggmencode}
\newcommand{\tableTL}{T_{\ggmencode_\plist}}

% Schwartz--Zippel game
\newcommand{\SZ}{\mathrm{SZ}}

% Source class
\newcommand{\sclass}{\mathbf{S}}


%----GAMES-----------------

% Games
\newcommand{\game}{\mathsf{Game}}
\newcommand{\collgame}{\mathsf{Coll}}
\newcommand{\pcind}{\text{}\quad}
\newcommand{\Bad}{\mathsf{Bad}} 
\newcommand{\Abort}{\mathsf{Abort}}

% Environment for game description
% Game transitions
\newcommand{\transitionarrow}{\rightsquigarrow}

% Font size for game hops in proofs
\newcommand{\gamesize}{\footnotesize}


%----CONSTRUCTIONS-----------------

%% Encryption schemes
% Syntax

% Security

\newcommand{\prelog}{\chi}
% Security
\newcommand{\UCE}{\mathrm{UCE}}
\newcommand{\uce}{\mathrm{uce}}
\newcommand{\HashO}{\textsc{Hash}}
\newcommand{\LabelO}{\textsc{Label}}
\newcommand{\LDD}{\mathrm{LDD}}
\newcommand{\ldd}{\mathrm{ldd}}
% Unpredictability
\newcommand{\Pred}{\mathrm{Pred}}
\newcommand{\pred}{\mathrm{pred}}
\newcommand{\ssup}{\mathrm{sup}}
\newcommand{\splt}{\mathrm{splt}}
\newcommand{\ssplt}{\mathrm{ssplt}}
\newcommand{\Splt}{\mathsf{Splt}}
\newcommand{\low}{\mathrm{low}}

%% RO-LDDs
\newcommand{\FSLDD}{\mathrm{FS} \dash \LDD}
\newcommand{\FSPred}{\mathrm{FS} \dash \Pred}
\newcommand{\MPLDD}{\MP \dash \LDD}
\newcommand{\MPPred}{\MP \dash \Pred}

%% Generic groups
% Syntax
\newcommand{\aigg}{\mathrm{AI} \dash \mathrm{GG}}
\newcommand{\bfgg}{\mathrm{BF} \dash \mathrm{GG}}
\newcommand{\ggmencode}{\tau}
\newcommand{\op}{\mathsf{op}}
\newcommand{\Decompose}{\textsc{Decomp}}

%% Concrete groups
\newcommand{\DLog}{\mathsf{DLog}}
% Security
\newcommand{\DUAII}{\mathrm{DUA \dash II}}
\newcommand{\duaii}{\mathrm{dua \dash ii}}
\newcommand{\DUA}{\mathrm{DUA}}
\newcommand{\dua}{\mathrm{dua}}
\newcommand{\DDHII}{\mathrm{DDH \dash II}}
\newcommand{\ddhii}{\mathrm{ddh \dash ii}}

%% PGGs
% Syntax
\newcommand{\encode}{\sigma}
\newcommand{\randfunc}{\rho}
% Security
\newcommand{\PGG}{\mathrm{PGG}}
\newcommand{\pgg}{\mathrm{pgg}}
% Unpredictability
\newcommand{\AlPred}{\mathrm{AlgPred}}
\newcommand{\alpred}{\mathrm{alg \dash pred}}
\newcommand{\MAlPred}{\mathrm{MAlgPred}}
\newcommand{\alg}{\mathrm{alg}}
\newcommand{\mask}{\mathrm{msk}}
\newcommand{\duber}{\mathrm{duber}}


%----NAMED CONSTRUCTIONS-----------------

\newcommand{\HtEl}{\mathsf{ElwH}}
\newcommand{\ShEl}{\mathsf{ShEl}}
\newcommand{\El}{\mathsf{El}}
\newcommand{\MEl}{\mathsf{MEl}}
\newcommand{\GOR}{\mathsf{GOR}}


\newcommand{\pad}{\mathsl{pad}}

\newcommand{\NGen}{\mathsf{NG}}
\newcommand{\Nonces}{\mathcal{N}}

\newcommand{\npars}{\sigma}
\newcommand{\NPred}{\gamefont{RP}}
\newcommand{\GenO}{\mathsc{Gen}}
%\newcommand{\EncO}{\mathsc{Enc}}
\newcommand{\HO}{\mathsc{H}}
\newcommand{\npred}{\mathrm{rp}}
\newcommand{\neScheme}{\schemefont{NE}}
\newcommand{\neVf}{\schemefont{NE.Vf}}
\newcommand{\neKg}{\neScheme.\schemefont{Kg}}
\newcommand{\neEnc}{\neScheme.\schemefont{Enc}}
\newcommand{\neDec}{\neScheme.\schemefont{Dec}}
\newcommand{\neSg}{\neScheme.\schemefont{Sg}}
\newcommand{\xk}{\varfont{xk}}
\newcommand{\vecxk}{\boldsymbol{xk}}
\newcommand{\vecN}{\boldsymbol{N}}
\newcommand{\vecStt}{\boldsymbol{st}}
\newcommand{\vecS}{\boldsymbol{S}}
\newcommand{\neNl}{\neScheme.\schemefont{nl}}
\newcommand{\seeds}{\boldsymbol{xk}}
\newcommand{\OpenO}{\mathsc{Open}}
\newcommand{\usr}{U}
\newcommand{\nind}{\ensuremath{\textrm{n-so-cpa}}}
\newcommand{\ncca}{\ensuremath{\textrm{n-so-cca}}}
\newcommand{\col}{\mathrm{cr}}
\newcommand{\GColl}{\mathrm{CR}}
\newcommand{\vecnpars}{\boldsymbol{\sigma}}
\newcommand{\neOneScheme}{\mathsf{NE1}}
\newcommand{\neOneKg}{\neOneScheme.\schemefont{Kg}}
\newcommand{\neOneEnc}{\neOneScheme.\schemefont{Enc}}
\newcommand{\neOneDec}{\neOneScheme.\schemefont{Dec}}
\newcommand{\neOneSg}{\neOneScheme.\schemefont{Sg}}
\newcommand{\neTwoScheme}{\mathsf{NE2}}
\newcommand{\neTwoKg}{\neTwoScheme.\schemefont{Kg}}
\newcommand{\neTwoEnc}{\neTwoScheme.\schemefont{Enc}}
\newcommand{\neTwoDec}{\neTwoScheme.\schemefont{Dec}}
\newcommand{\neTwoSg}{\neTwoScheme.\schemefont{Sg}}


\newcommand{\HNCCA}{\ensuremath{\textrm{HN-SO-CCA}}}
\newcommand{\HNCPA}{\ensuremath{\textrm{HN-SO-CPA}}}
\newcommand{\hncpa}{\ensuremath{\textrm{hn-so-cpa}}}
\newcommand{\hncca}{\ensuremath{\textrm{hn-so-cca}}}
\newcommand{\HNCPAREAL}{\ensuremath{\textrm{HN-CPA-REAL}}}
\newcommand{\HNCCAREAL}{\ensuremath{\textrm{HN-CCA-REAL}}}
\newcommand{\HNCPAIDEAL}{\ensuremath{\textrm{HN-CPA-IDEAL}}}
\newcommand{\HNCCAIDEAL}{\ensuremath{\textrm{HN-CCA-IDEAL}}}

\newcommand{\neThreeScheme}{\mathsf{NE1}}
\newcommand{\neThreeKg}{\neThreeScheme.\schemefont{Kg}}
\newcommand{\neThreeEnc}{\neThreeScheme.\schemefont{Enc}}
\newcommand{\neThreeDec}{\neThreeScheme.\schemefont{Dec}}
\newcommand{\neThreeSg}{\neThreeScheme.\schemefont{Sg}}
\newcommand{\neFourScheme}{\mathsf{NE2}}
\newcommand{\neFourKg}{\neFourScheme.\schemefont{Kg}}
\newcommand{\neFourEnc}{\neFourScheme.\schemefont{Enc}}
\newcommand{\neFourDec}{\neFourScheme.\schemefont{Dec}}
\newcommand{\neFourSg}{\neFourScheme.\schemefont{Sg}}

\newcommand{\neBadScheme}{\mathsf{NE_{bad}}}
\newcommand{\neBadKg}{\neBadScheme.\schemefont{Kg}}
\newcommand{\neBadEnc}{\neBadScheme.\schemefont{Enc}}
\newcommand{\neBadDec}{\neBadScheme.\schemefont{Dec}}
\newcommand{\neBadSg}{\neBadScheme.\schemefont{Sg}}

\newcommand{\NtD}{\mathsf{NtD}}

\newcommand{\neeScheme}{\mathsf{\overline{NE}}}
\newcommand{\neeKg}{\neeScheme.\schemefont{Kg}}
\newcommand{\neeEnc}{\neeScheme.\schemefont{Enc}}
\newcommand{\neeDec}{\neeScheme.\schemefont{Dec}}
\newcommand{\neeSg}{\neeScheme.\schemefont{Sg}}

\newcommand{\Hbar}{\overline{H}}
\newcommand{\advAbar}{\overline{\advA}}
\newcommand{\veccbar}{\overline{\vecc}}

\newcommand{\clear}{\mathsf{\zeta\mbox{-}clear}}


\newcommand{\padOAEPb}{
\psset{arrowsize=0.2}
\begin{psmatrix}[colsep=0.5cm,rowsep=0.1cm]
&m  & &  \\[-1.5ex]
&[name=m]&&[name=r]&\\[1.0ex]
&[name=p1]\oplus&[name=H1]\psframebox{G}&[name=x1]\\[1.0ex]
&[name=x2]&[name=H2]\psframebox{H}&[name=p2]\oplus\\[1.0ex]
&[name=s]&&[name=t]\\[-1.5ex]
&s\in\{0,1\}^\ms &&t \in \{0,1\}^\rs
\end{psmatrix}
\ncline{<-}{m}{s}
\ncline{-}{x1}{t}
\ncline{-}{p1}{H1}
\ncline{-}{x1}{H1}
\ncline{-}{p2}{H2}
\ncline{-}{x2}{H2}
}





\newcommand{\PADScheme}{\mathsf{PAD}}


\newcommand{\PKscheme}{{\cal AE}}

\renewcommand{\P}{\mathrm{Pr}}
\newcommand{\probExp}[2]{\P\bigl[\, #1 \given #2\, \bigr]}

\newcommand{\OAEP}{\mathsf{OAEP}}

\newcommand{\rs}{\rho}
\newcommand{\ms}{\mu}

\newcommand{\PAD}{\Pi}
\newcommand{\padi}{\hat{\pi}}

\renewcommand{\b}[1]{\text{\boldmath$#1$\unboldmath}}

\newcommand{\rleak}{s}
\newcommand{\eps}{\varepsilon}

\newcommand{\LTDP}{\mathsf{LTDP}}

\newcommand{\COMEXT}{\ensuremath{\textrm{COM-EXT}}}
\newcommand{\comext}{\ensuremath{\textrm{com-ext}}}


%%% Local Variables: 
%%% mode: latex
%%% TeX-master: "main-full"
%%% End: 

\newcommand{\vechint}{\mathbf{hint}}
\newcommand{\hint}{\mathsf{hint}}


\newcommand{\cd}{\algfont{CD}}

\newcommand{\coin}{\mathsf{coin}}

\newcommand{\timet}{\algfont{v}}

%\newcommand{\FDom}{\mathsf{FDom}}
%\newcommand{\FRng}{\mathsf{FRng}}
\newcommand{\CIHGame}{\ensuremath{\textrm{2HCF}\mbox{-}\textrm{AI}}}

% To do later
% - enlarge brackets also in inline text
% - macro ( )
% - interval for []
% - add ~ before math
% - penalize inline math breaking?
% - decide on \cdot policy
% - make font of games in proofs smaller
% - check paragraph indentation
% - + at end of line if it doesn't fit completely
% - Add all (\secpar) in transitions?
% - Right way of embedding text into math (for game "sets" bad)
% - remember to add (and open) bookmarks

% - make lcnamecref work in the appendix, so that it shows section there

% Authnotes
\newcommand{\adam}[1]{{\color{blue}[Adam: #1]}}
\newcommand{\pooya}[1]{{\color{blue}[Pooya: #1]}}